\SourceLeaf{502}die die relative Bewegung des Körpers gegen die bewegte Erdoberfläche in sich fassen, aus den ersten Fundamentalsätzen der Dynamik, und hier fand ich zu meiner Verwunderung

1) die Abweichung nach Süden wiederum $0$ oder ganz unvermerklich:

2) die Abweichung nach Osten nur $\tfrac23$ von dem, was Dr.\ \textsc{Olbers} gefunden hat. Nemlich in Dr.\ \textsc{Olbers} Zeichen, wenn man den Widerstand der Luft vernachlässigt,
\[
=\frac{\tfrac43\pi\cos\psi.at}{86164}
\]
oder wenn man ihn mit in Betrachtung zieht, nach einer hier zureichenden Näherung
\[
=\frac{\tfrac43\pi\cos\psi.t}{86164}(\tfrac32a'-\tfrac12a)
\]
wo $a$ die Höhe ist, durch die der Körper in der Zeit $t$ im leeren Raume fallen würde, also $=\tfrac12g'tt$ [wo ferner $\psi$ die Polhöhe des Beobachtungsortes und $a'$ die wirkliche Fallhöhe bezeichnet].

Hienach finde ich für Ihre Versuche, indem ich die Pendellänge für Hamburg $=440.75$ Linien (woraus $g'$ fast eben so kommt, wie Dr.\ \textsc{Olbers} es annimmt), \textit{die Abweichung nach Osten} $3.951$ \textit{pariser Linien}; welches sehr genau mit Ihren Versuchen übereinstimmte, — da hingegen die Abweichung nach Süden nicht zu meinen Resultaten passt.

\begin{center}\resizebox{!}{58mm}{\begin{tikzpicture}[x=1bp,y=-1bp,line cap=round,line join=round]
\path[use as bounding box] (0,0) rectangle (395,752);
\begin{pgfinterruptboundingbox}
\draw[line width=2.3bp] (75,215) .. controls (140,109) and (253,62) .. (326,56);
\draw[line width=1.8bp] (75,215) -- (327,218);
\draw[line width=1.7bp] (327,56) -- (327,85);
\draw[line width=1.7bp] (327,106) -- (326,257);
\draw[line width=1.7bp] (326,270) -- (324,316);
\draw[line width=1.7bp] (324,326) -- (323,411);
\draw[line width=1.7bp] (323,438) -- (317,684);
\node[inner sep=0pt,outer sep=0pt] at (315,37) {\scalebox{3.8}{\ensuremath{m}}};
\node[inner sep=0pt,outer sep=0pt] at (45,201) {\scalebox{3.8}{\ensuremath{M}}};
\node[inner sep=0pt,outer sep=0pt] at (347,207) {\scalebox{3.8}{\ensuremath{R}}};
\node[inner sep=0pt,outer sep=0pt] at (316,708) {\scalebox{3.8}{\ensuremath{C}}};
\end{pgfinterruptboundingbox}
\end{tikzpicture}
}\end{center}
Diese Verschiedenheit in Ansehung der Abweichung nach Osten — veranlasste mich, Dr.\ \textsc{Olbers} Schlüsse darüber aufmerksamer durchzugehen, und die Ursache davon nachzuspüren. Wie mir scheint, liegt sie darin, dass Dr.\ \textsc{Olbers} die wirkliche Bewegung des Körpers gegen Osten \textit{bloss} aus seiner tangentiellen ursprünglichen Geschwindigkeit ableitet, und von der daraus entspringenden Bewegung die gleichzeitige Bewegung des Fusses des Thurms abzieht, um die scheinbare Bewegung nach Osten zu haben. — Allein wenn die Fläche des Papiers die obige Ebene $\mathrm{A}$ vorstellt, $\mathrm{C}$ den Mittelpunkt der Erde, $\mathrm{mM}$ die wirkliche Bewegung des Körpers: so darf man, meiner Meinung nach, nicht ausser Acht lassen, dass selbst die Anziehung nach $\mathrm{C}$ während die Bewegung nicht mit $\mathrm{mC}$ parallel ist, und eben da\SourceLeaf{503}her die Geschwindigkeit nach Osten wirklich vermindert wird, daher der Körper, wenn er in $\mathrm{M}$ anlangt, nicht so weit nach Osten gekommen ist, als er mit der ursprünglichen Geschwindigkeit gekommen sein würde. Nach darüber geführter Rechnung finde ich auch, dass durch diese Betrachtung die scheinbare Bewegung nach Osten wirklich um den dritten Theil vermindert wird.

\Needspace{7\baselineskip}
\begin{center}{\large\itshape Brief von Gauss an Benzenberg.}\end{center}
\begin{flushright}Braunschweig 1803. März 8.\end{flushright}
—\quad—\quad—\quad—\quad An unsern Freund \textsc{Olbers} habe ich vor acht Tagen einen kleinen Aufsatz über die Abweichung fallender Körper eingesandt. Heute erhalte ich darauf die Antwort:

1) die Abweichung nach Osten sei nur $\tfrac23$ von der, die er berechnet hätte:

2) dass er meinen Schlüssen, dass die Abweichung nach Süden $=0$ sei, nichts entgegenzusetzen habe. aber zu wissen wünsche, \textit{worin} eigentlich \textit{sein} Raisonnement fehlerhaft sei.

Ich bemerke hiebei noch folgendes:

Vorausgesetzt, dass meine Schlüsse in Ansehung der Abweichung nach Süden gewiss sind, so scheint mir der Grund von der von Dr.\ \textsc{Olbers} herausgebrachten Abweichung noch immer darin zu liegen, dass er voraussetzt, der Körper \textit{bleibe} auch bei widerstehender Luft in der auf den Meridian senkrechten, und durch den Mittelpunkt der Erde gehenden Ebene. Es scheint mir, dass diese Voraussetzung nothwendig gerechtfertigt werden müsse, aber ich zweifle, ob sie sich rechtfertigen lasse. Die kegelförmige Bewegung der Luft macht, dass die Lufttheile, worin der Körper ist, sobald die Erde aus ihrer ersten Lage gekommen, in einem Winkel durch jene Ebene gehn, den man nicht vernachlässigen darf, und wodurch es geschieht, dass der Körper, dem die Luft nicht in der Richtung dieser Ebene widersteht, aus der Ebene gegen Norden heraustritt: und ich bin noch immer der Meinung, dass sie aus der Verzögerung dadurch vollkommen compensirt wird. Es ist mir auch wahrscheinlich, dass \textsc{Guglielmini} eben dies hat sagen wollen, und dass er nur deswegen \textsc{Olbers} Beifall nicht erhalten hat, weil er sich nicht bestimmt genug erklärt. Ich hoffe indess zuversichtlich, dass entweder ich mit Dr.\ \textsc{Olbers}, oder Dr.\ \textsc{Olbers} mit mir vollkommen zu einerlei Überzeugung kommen werden. —\quad—\quad—

\SourceLeaf{504}\Needspace{8\baselineskip}
\begin{center}{\large\itshape Fundamentalgleichungen für die Bewegung schwerer Körper\par auf der rotirenden Erde.}\end{center}
Die Lage eines Punkts wird auf eine doppelte Art bestimmt.

\textit{Erstens} durch seine senkrechten Abstände $X,Y,Z$, von drei auf einander senkrechten \textit{festen} Ebnen. Den gemeinschaftlichen Durchschnittspunkt dieser Ebnen, $C$, setzen wir in einen beliebigen Punkt der Erdaxe; die Ebene der $Z$ legen wir dem Aequator parallel; die Ebene der $Y$ in denjenigen Meridian, worin sich der anfängliche Ort des Körpers befindet; endlich die Ebene der $X$ in den auf den vorigen senkrechten Meridian. Die $Z$ sind positiv auf der Nordseite; die $X$ auf der Seite des anfänglichen Orts des Körpers, die $Y$ auf derjenigen Seite, wohin dieser anfängliche Ort durch die Rotation geführt wird.

\textit{Zweitens} durch die senkrechten Abstände $x,y,z$, von drei auf einander senkrechten \textit{beweglichen} d.\,i. gegen die Erde ruhenden und mit ihr rotirenden Ebnen. Am schicklichsten setzen wir den gemeinschaftlichen Durchschnittspunkt derselben in den anfänglichen Ort des Körpers. Die Ebene der $z$ setzen wir senkrecht auf die scheinbare Richtung der Schwere; die der $y$ in den Meridian: dadurch ist die auf beide senkrechte der $x$ von selbst bestimmt; Pole dieser drei Ebnen sind also resp.\ das scheinbare Zenith, der Ostpunkt, der Südpunkt, und \textit{diese} Pole sollen zugleich diejenigen Seiten der Ebnen bezeichnen, wo die Abstände $z,y,x$ positiv genommen werden.

Es sei jetzt für den Punkt $C$, $x=a$, $(y=0)$, $z=-c$; ferner die (scheinbare, nördliche) Polhöhe des Beobachtungsorts $\varphi$, und der Winkel, um den sich die Erde nach der Zeit $t$ gegen Osten bewegt hat, $\theta$. Unter diesen Voraussetzungen ergeben sich leicht folgende Gleichungen:
\[
\left.\begin{aligned}
x&=X\sin\varphi\cos\theta+Y\sin\varphi\sin\theta-Z\cos\varphi+a\\
y&=-X\sin\theta+Y\cos\theta\\
z&=X\cos\varphi\cos\theta+Y\cos\varphi\sin\theta+Z\sin\varphi-c
\end{aligned}\right\}\quad.\quad.\quad.\quad.\quad.\quad[1]
\]
\[
\left.\begin{aligned}
X&=(x-a)\sin\varphi\cos\theta-y\sin\theta+(z+c)\cos\varphi\cos\theta\\
Y&=(x-a)\sin\varphi\sin\theta+y\cos\theta+(z+c)\cos\varphi\sin\theta\\
Z&=-(x-a)\cos\varphi+(z+c)\sin\varphi
\end{aligned}\right\}\quad.\quad.\quad.\quad[2]
\]

\SourceLeaf{505}Die Coordinaten $X,Y,Z$ lassen sich einerseits als Functionen von $t$ allein, anderseits aber auch als Functionen der vier veränderlichen Grössen $\theta,x,y,z$ betrachten, und haben also in letzterer Hinsicht vier partielle Differentiale. Es ist demnach
\[
\begin{aligned}
\mathrm{d}X&=(\tfrac{\mathrm{d}X}{\mathrm{d}t})\mathrm{d}t=(\tfrac{\mathrm{d}X}{\mathrm{d}\theta})\mathrm{d}\theta+(\tfrac{\mathrm{d}X}{\mathrm{d}x})\mathrm{d}x+(\tfrac{\mathrm{d}X}{\mathrm{d}y})\mathrm{d}y+(\tfrac{\mathrm{d}X}{\mathrm{d}z})\mathrm{d}z\\
\mathrm{d}Y&=\mathrm{etc.}
\end{aligned}
\]
Die Geschwindigkeit des Körpers zerlegt sich, wie seine Bewegung, in drei partielle auf die Ebnen der $X,Y,Z$ senkrechte Geschwindigkeiten, die mithin $(\tfrac{\mathrm{d}X}{\mathrm{d}t})$, $(\tfrac{\mathrm{d}Y}{\mathrm{d}t})$, $(\tfrac{\mathrm{d}Z}{\mathrm{d}t})$ sind. Die Geschwindigkeiten des Luftelements hingegen, in welchem er sich jedesmal befindet, in Beziehung auf dieselben Ebnen sind offenbar $(\tfrac{\mathrm{d}X}{\mathrm{d}\theta})\tfrac{\mathrm{d}\theta}{\mathrm{d}t}$, $(\tfrac{\mathrm{d}Y}{\mathrm{d}\theta})\tfrac{\mathrm{d}\theta}{\mathrm{d}t}$, $(\tfrac{\mathrm{d}Z}{\mathrm{d}\theta})\tfrac{\mathrm{d}\theta}{\mathrm{d}t}$. Folglich die \textit{relativen} Geschwindigkeiten des Körpers nach diesen drei Richtungen
\[
\begin{aligned}
(\tfrac{\mathrm{d}X}{\mathrm{d}x})\tfrac{\mathrm{d}x}{\mathrm{d}t}+(\tfrac{\mathrm{d}X}{\mathrm{d}y})\tfrac{\mathrm{d}y}{\mathrm{d}t}+(\tfrac{\mathrm{d}X}{\mathrm{d}z})\tfrac{\mathrm{d}z}{\mathrm{d}t}&=\xi=\sin\varphi\cos\theta\tfrac{\mathrm{d}x}{\mathrm{d}t}-\sin\theta\tfrac{\mathrm{d}y}{\mathrm{d}t}+\cos\varphi\cos\theta\tfrac{\mathrm{d}z}{\mathrm{d}t}\\
(\tfrac{\mathrm{d}Y}{\mathrm{d}x})\tfrac{\mathrm{d}x}{\mathrm{d}t}+(\tfrac{\mathrm{d}Y}{\mathrm{d}y})\tfrac{\mathrm{d}y}{\mathrm{d}t}+(\tfrac{\mathrm{d}Y}{\mathrm{d}z})\tfrac{\mathrm{d}z}{\mathrm{d}t}&=\eta=\sin\varphi\sin\theta\tfrac{\mathrm{d}x}{\mathrm{d}t}+\cos\theta\tfrac{\mathrm{d}y}{\mathrm{d}t}+\cos\varphi\sin\theta\tfrac{\mathrm{d}z}{\mathrm{d}t}\\
(\tfrac{\mathrm{d}Z}{\mathrm{d}x})\tfrac{\mathrm{d}x}{\mathrm{d}t}+(\tfrac{\mathrm{d}Z}{\mathrm{d}y})\tfrac{\mathrm{d}y}{\mathrm{d}t}+(\tfrac{\mathrm{d}Z}{\mathrm{d}z})\tfrac{\mathrm{d}z}{\mathrm{d}t}&=\zeta=-\cos\varphi\tfrac{\mathrm{d}x}{\mathrm{d}t}+\sin\varphi\tfrac{\mathrm{d}z}{\mathrm{d}t}
\end{aligned}
\]
Die totale relative Geschwindigkeit ist folglich $=\BareRoot(\xi\xi+\eta\eta+\zeta\zeta)=u$, welches, wie die Entwickelung aus obigen Werthen leicht zeigt, $=\BareRoot(\tfrac{\mathrm{d}x^2}{\mathrm{d}t^2}+\tfrac{\mathrm{d}y^2}{\mathrm{d}t^2}+\tfrac{\mathrm{d}z^2}{\mathrm{d}t^2})$ wird. Der Widerstand der Luft ist dem Quadrate davon proportional, wir setzen ihn daher $=Muu$, und zerlegen ihn nach obigen drei Richtungen in $Mu\xi$, $Mu\eta$, $Mu\zeta$.

Wir sehen hier die Erde als ein Revolutions-Sphäroid an; die Richtung der Schwere geht daher durch die Erdaxe. Der Punkt, wo diese schneidet, liege um $q$ über $C$, oder es sei für denselben $Z=q$.

Setzt man nun ferner die Stärke der Gravitation $=p$ und
\[
XX+YY+(Z-q)^2=rr
\]
so ist nach den Grundsätzen der Dynamik
\SourceLeaf{506}\[
\left.\begin{aligned}
0&=\frac{\mathrm{dd}X}{\mathrm{d}t^2}+\frac{pX}{r}+Mu\xi\\
0&=\frac{\mathrm{dd}Y}{\mathrm{d}t^2}+\frac{pY}{r}+Mu\eta\\
0&=\frac{\mathrm{dd}Z}{\mathrm{d}t^2}+\frac{p(Z-q)}{r}+Mu\zeta
\end{aligned}\right\}\quad.\quad.\quad.\quad.\quad.\quad.\quad.\quad[3]
\]
Aus obigen Werthen von $X,Y,Z$ in [2] findet man, wenn man für $\tfrac{\mathrm{d}\theta}{\mathrm{d}t}$, welches beständig ist, $n$ schreibt, folgende Gleichungen:
\[
\begin{aligned}
\frac{\mathrm{dd}X}{\mathrm{d}t^2}={}&\sin\varphi\cos\theta\frac{\mathrm{dd}x}{\mathrm{d}t^2}-\sin\theta\frac{\mathrm{dd}y}{\mathrm{d}t^2}+\cos\varphi\cos\theta\frac{\mathrm{dd}z}{\mathrm{d}t^2}\\
&-2n\sin\varphi\sin\theta\frac{\mathrm{d}x}{\mathrm{d}t}-2n\cos\theta\frac{\mathrm{d}y}{\mathrm{d}t}-2n\cos\varphi\sin\theta\frac{\mathrm{d}z}{\mathrm{d}t}-nnX\\[5pt]
\frac{\mathrm{dd}Y}{\mathrm{d}t^2}={}&\sin\varphi\sin\theta\frac{\mathrm{dd}x}{\mathrm{d}t^2}+\cos\theta\frac{\mathrm{dd}y}{\mathrm{d}t^2}+\cos\varphi\sin\theta\frac{\mathrm{dd}z}{\mathrm{d}t^2}\\
&+2n\sin\varphi\cos\theta\frac{\mathrm{d}x}{\mathrm{d}t}-2n\sin\theta\frac{\mathrm{d}y}{\mathrm{d}t}+2n\cos\varphi\cos\theta\frac{\mathrm{d}z}{\mathrm{d}t}-nnY\\[5pt]
\frac{\mathrm{dd}Z}{\mathrm{d}t^2}={}&-\cos\varphi\frac{\mathrm{dd}x}{\mathrm{d}t^2}+\sin\varphi\frac{\mathrm{dd}z}{\mathrm{d}t^2}
\end{aligned}\quad[4]
\]
Multiplicirt man die drei Gleichungen [3] resp.\ mit $\sin\varphi\cos\theta$, $\sin\varphi\sin\theta$, $-\cos\varphi$ und addirt die Producte; multiplicirt man zweitens eben diese Gleichungen mit $-\sin\theta$, $\cos\theta$, $0$; und drittens mit $\cos\varphi\cos\theta$, $\cos\varphi\sin\theta$, $\sin\varphi$, und addirt beidemale die Producte: so erhält man, nachdem man statt $\tfrac{\mathrm{dd}X}{\mathrm{d}t^2}$, $\tfrac{\mathrm{dd}Y}{\mathrm{d}t^2}$, $\tfrac{\mathrm{dd}Z}{\mathrm{d}t^2}$ ihre Werthe aus [4], statt $X,Y,Z$ die aus [2], und statt $\xi$, $\eta$, $\zeta$ die ihrigen substituirt hat, folgende drei neue:
\[
\begin{aligned}
0&=\tfrac{\mathrm{dd}x}{\mathrm{d}t^2}-2n\sin\varphi\tfrac{\mathrm{d}y}{\mathrm{d}t}+(x-a)(\tfrac pr-nn)+\cos\varphi(\tfrac{pq}{r}-nnZ)+Mu\tfrac{\mathrm{d}x}{\mathrm{d}t}\\
0&=\tfrac{\mathrm{dd}y}{\mathrm{d}t^2}+2n\sin\varphi\tfrac{\mathrm{d}x}{\mathrm{d}t}+2n\cos\varphi\tfrac{\mathrm{d}z}{\mathrm{d}t}+y(\tfrac pr-nn)+Mu\tfrac{\mathrm{d}y}{\mathrm{d}t}\\
0&=\tfrac{\mathrm{dd}z}{\mathrm{d}t^2}-2n\cos\varphi\tfrac{\mathrm{d}y}{\mathrm{d}t}+(z+c)(\tfrac pr-nn)-\sin\varphi(\tfrac{pq}{r}-nnZ)+Mu\tfrac{\mathrm{d}z}{\mathrm{d}t}
\end{aligned}
\]
Ist also der Körper gegen die Erde in Ruhe, oder $\mathrm{d}x=\mathrm{d}y=\mathrm{d}z=0$, so scheint er senkrecht auf die Ebnen der $x,y,z$ von den Kräften
\SourceLeaf{507}\[
\begin{gathered}
(x-a)(\tfrac pr-nn)+\cos\varphi(\tfrac{pq}{r}-nnZ)\\
y(\tfrac pr-nn)\\
(z+c)(\tfrac pr-nn)-\sin\varphi(\tfrac{pq}{r}-nnZ)
\end{gathered}
\]
sollicitirt zu werden. Ein schon in Bewegung begriffener Körper hingegen wird anders afficirt. Denn ausser dem Widerstande der Luft, der den Körper nach diesen Richtungen wie Kräfte, deren Maass $Mu\tfrac{\mathrm{d}x}{\mathrm{d}t}$, $Mu\tfrac{\mathrm{d}y}{\mathrm{d}t}$, $Mu\tfrac{\mathrm{d}z}{\mathrm{d}t}$ ist, treibt und folglich auf der rotirenden Erde völlig eben so wirkt, als er auf der ruhenden wirken würde, kommen nach jenen Richtungen noch die drei Kräfte
\[
-2n\sin\varphi\frac{\mathrm{d}y}{\mathrm{d}t},\quad 2n\sin\varphi\frac{\mathrm{d}x}{\mathrm{d}t}+2n\cos\varphi\frac{\mathrm{d}z}{\mathrm{d}t},\quad -2n\cos\varphi\frac{\mathrm{d}y}{\mathrm{d}t}
\]
hinzu, und diese sind es allein, wodurch die Rotation der Erde an fallenden Körpern sichtbar wird. Die bisherigen Schlüsse und Folgerungen sind streng und allgemein richtig.

Bei \textit{Versuchen}, die in dieser Hinsicht angestellt werden, geschieht allemal die Bewegung des Körpers in einem so kleinen Raume, dass man die Stärke der auf ruhende Körper wirkenden scheinbaren Schwere innerhalb desselben, als unveränderlich $=g$, und ihre Richtung als immer parallel, also senkrecht auf die Ebene der $z$ annehmen kann. Es wird also ohne Bedenken erlaubt sein, statt der obigen drei Grössen
\[
\begin{gathered}
(x-a)(\tfrac pr-nn)+\cos\varphi(\tfrac{pq}{r}-nnZ)\\
y(\tfrac pr-nn)\\
(z+c)(\tfrac pr-nn)-\sin\varphi(\tfrac{pq}{r}-nnZ)
\end{gathered}
\]
respective $0,0,g$ zu substituiren. Dadurch werden die drei Fundamentalgleichungen
\[
\begin{aligned}
0&=\frac{\mathrm{dd}x}{\mathrm{d}t^2}-2n\sin\varphi\frac{\mathrm{d}y}{\mathrm{d}t}+Mu\frac{\mathrm{d}x}{\mathrm{d}t}\\
0&=\frac{\mathrm{dd}y}{\mathrm{d}t^2}+2n\sin\varphi\frac{\mathrm{d}x}{\mathrm{d}t}+2n\cos\varphi\frac{\mathrm{d}z}{\mathrm{d}t}+Mu\frac{\mathrm{d}y}{\mathrm{d}t}\\
0&=\frac{\mathrm{dd}z}{\mathrm{d}t^2}-2n\cos\varphi\frac{\mathrm{d}y}{\mathrm{d}t}+g+Mu\frac{\mathrm{d}z}{\mathrm{d}t}
\end{aligned}
\]
\SourceLeaf{508}Die Integration dieser Gleichungen ist leicht, wenn man den Widerstand der Luft vernachlässigt, oder $M=0$ setzt. Man findet nemlich
\[
\begin{aligned}
x&=\mathfrak{A}-\mathfrak{D}\cos\varphi.t+\tfrac{1}{2n}\mathfrak{E}\sin\varphi\cos(2nt+\mathfrak{F})+\tfrac12\sin\varphi\cos\varphi.gtt\\
y&=\mathfrak{B}-\tfrac{1}{2n}\mathfrak{E}\sin(2nt+\mathfrak{F})+\tfrac{1}{2n}\cos\varphi.gt\\
z&=\mathfrak{C}+\mathfrak{D}\sin\varphi.t+\tfrac{1}{2n}\mathfrak{E}\cos\varphi\cos(2nt+\mathfrak{F})-\tfrac12\sin\varphi^2.gtt
\end{aligned}
\]
Auch ist es leicht, folgende Werthe der arbiträren Grössen zu entwickeln, wenn man voraussetzt, dass der Körper anfänglich gar keine scheinbare Geschwindigkeit hat:
\[
\begin{array}{lll}
\mathfrak{A}=-\tfrac{g}{4nn}\cos\varphi\sin\varphi,&\mathfrak{B}=0,&\mathfrak{C}=-\tfrac{g}{4nn}\cos\varphi^2\\
\mathfrak{D}=0,&\mathfrak{E}=\tfrac{g}{2n}\cos\varphi,&\mathfrak{F}=0
\end{array}
\]
Also
\[
\begin{aligned}
x&=\tfrac{g}{2n}\cos\varphi\sin\varphi(ntt-\tfrac{1}{2n}+\tfrac{1}{2n}\cos2nt)\\
y&=\tfrac{g}{2n}\cos\varphi(t-\tfrac{1}{2n}\sin2nt)\\
z&=-\tfrac12gtt+\tfrac{g}{2n}\cos\varphi^2(ntt-\tfrac{1}{2n}+\tfrac{1}{2n}\cos2nt)
\end{aligned}
\]
Diese Integration ist freilich nicht \textit{allgemein} zulässig, da obige Voraussetzung nur in so fern erlaubt ist, als der Körper sich von seinem anfänglichen scheinbaren Orte nicht weit entfernt. Für diesen Fall aber können wir die trigonometrischen Functionen in Reihen auflösen, und so wird
\[
\begin{aligned}
x&=\tfrac16\cos\varphi\sin\varphi.gnnt^4\quad.\quad.\quad.\\
y&=\tfrac13\cos\varphi.gnt^3\quad.\quad.\quad.\\
z&=-\tfrac12gtt+\tfrac16\cos\varphi^2.gnnt^4\quad.\quad.\quad.
\end{aligned}
\]
Da die Zeit des Falls nur wenige Secunden, also $nt$ höchstens einige Raumminuten beträgt, und (weil Radius $=1$) $1'=\tfrac{1}{3438}$, so wird $x$ und der zweite Theil von $z$ ganz unmerklich, also $y=-\tfrac23z\cos\varphi.nt$. Bei Dr.\ \textsc{Benzenbergs} Versuche im Michaelisthurme war $z=-235$ Fuss, $\varphi=53^0\,33'$, $t=4''$ Sonnenzeit, also $nt=\tfrac{366}{365}$ Raumminuten. Hieraus wird $y=3{,}91$ Linien.

\SourceLeaf{509}Wenn man bei der Integration obiger Gleichungen den Widerstand der Luft mit in Betrachtung ziehen will, so wird man sich mit Näherungen begnügen müssen; die Entwickelung der Werthe von $x,y,z$ in Reihen nach den Potenzen von $n$ und $M$ ist alsdann sehr leicht. Das höchste Glied von $x$ wird wie vorhin $=\tfrac16\cos\varphi\sin\varphi.gnnt^4$, und ist also von gar keiner Bedeutung; für $y$ und $z$ findet man mit Vernachlässigung der Quadrate und höhern Potenzen von $n$ und $M$ folgende Werthe:
\[
\begin{aligned}
y&=\tfrac13\cos\varphi.gnt^3-\tfrac{1}{12}\cos\varphi.Mggnt^5\\
z&=-\tfrac12gtt+\tfrac{1}{12}Mggt^4.
\end{aligned}
\]
Setzen wir also $-z$, den wirklichen Fall, $=f$; $\tfrac12gtt$ oder den Fall im luftleeren Raume $=f+\delta$, so ist
\[
y=\tfrac23\cos\varphi.nt(f+\delta)-\cos\varphi.nt\delta=\tfrac23\cos\varphi.nt(f-\tfrac12\delta)
\]
Für die Versuche in St.\ Michael, wo $f+\delta=241{,}47$ Fuss war, erhalten wir daher die Abweichung nach Osten $y=3{,}86$ Linien.

\clearpage
\fancyhead[L]{\small ACHROMATISCHE DOPPELOBJECTIVE.}
\SourceLeaf{510}\pdfbookmark[0]{Über die achromatischen Doppelobjective}{achromatische_doppelobjective}
\begin{center}
{\Large ÜBER DIE ACHROMATISCHEN DOPPELOBJECTIVE\par}
\vspace{6mm}
{\large BESONDERS IN RÜCKSICHT\par}
\vspace{4mm}
{\large DER VOLLKOMMENEN AUFHEBUNG DER FARBENZERSTREUUNG.\par}
\vspace{7mm}\hrule\vspace{2mm}
Zeitschrift für Astronomie und verwandte Wissenschaften\par
herausgegeben von B.\ von \textsc{Lindenau} und \textsc{Bohnenberger}. Bd.\ IV. N.\ XXX. 1817. December.
\vspace{2mm}\hrule\vspace{5mm}
\end{center}
Der schöne Aufsatz des Hrn.\ Prof.\ \textsc{Bohnenberger} über die achromatischen Objective im ersten Bande dieser Zeitschrift hat das Verdienst, einen für diese Theorie wichtigen Umstand zuerst zur Sprache gebracht zu haben. Ich bin dadurch veranlasst, einige frühere Untersuchungen wieder vorzunehmen und weiter zu entwickeln, deren Resultate ich hier mittheilen werde.

Man begnügte sich bisher bei den Doppelobjectiven, die Farbenzerstreuung für die der Axe unendlich nahen Strahlen, und die Abweichung wegen der Kugelgestalt für die Strahlen von mittlerer Brechbarkeit zu heben, wobei also für die Randstrahlen noch eine kleine Farbenzerstreuung zurückbleiben kann. Bei dieser Einrichtung ist die Berechnung des achromatischen Objectivs eine unbestimmte Aufgabe, d.\,i., zu jeder Kronglaslinse von positiver Brennweite, wie auch immer das Verhältniss der Halbmesser der Flächen sein mag, lässt sich eine Flintglaslinse berechnen, die mit jener vereinigt ein in obiger Bedeutung achromatisches Objectiv gibt. So viel ich weiss, haben bisher alle Optiker beide Flächen der Kronglaslinse convex angenommen: allein für das Verhältniss der beiden Halbmesser haben die Theoretiker sehr verschiedene Werthe in Vorschlag gebracht, je nachdem sie von diesem oder jenem Princip ausgingen. Will man mit \textsc{Euler} die Abweichung wegen der Gestalt bei der Kronglaslinse zu einem Kleinsten machen, so müssen die Halbmesser ungefähr in dem Verhältniss von
\SourceLeaf{511}1 zu 7 stehen; sie müssen einander gleich sein, wenn man, wie \textsc{Klügel} in der analytischen Dioptrik, die möglich kleinsten Krümmungen zu haben wünscht; sollen die Brechungen selbst die möglich kleinsten werden, wie derselbe Schriftsteller in einer spätern Abhandlung sich vorsetzt, so müssen diese Brechungen einander gleich sein, und die Halbmesser nahe in dem Verhältniss von 1 zu 3 stehen. Es scheint nicht, dass alle diese verschiedenen Vorschläge hinlänglich motivirt sind. \textsc{Klügels} Augenmerk ist besonders die Abweichung wegen der Kugelgestalt gewesen, welche für \textit{alle} Strahlen in mathematischer Schärfe zu heben bekanntlich unmöglich ist: bei \textsc{Eulers} Behandlung dieser Rechnungen ist diese Abweichung eigentlich nur für die der Axe nächsten Strahlen gehoben, und es bleibt eine sehr nahe dem Biquadrat des Abstandes von der Axe proportionale, also für die Randstrahlen am meisten merkliche Abweichung zurück; oder wenn man mit \textsc{Klügel} die Rechnung so führt, dass die Abweichung für die Randstrahlen verschwindet, so kommt sie wieder bei den Zwischenstrahlen zum Vorschein, am merklichsten bei denen, deren Entfernung nahe $\tfrac{7}{10}$ von dem Halbmesser der Öffnung ist. Diese unvermeidlich übrigbleibende Abweichung wegen der Gestalt so unschädlich wie möglich zu machen, war \textsc{Klügels} Absicht bei der Wahl des Verhältnisses der beiden ersten Halbmesser: es erhellt jedoch nicht klar genug weder, dass wirklich dieser Zweck bei dem gewählten Verhältniss am allerbesten erreicht werde, noch, dass dieser Zweck wichtig genug sei, um ihn vorzugsweise allein zur Grundlage der Bestimmung dieses Verhältnisses zu machen. Finden nemlich noch \textit{andere} Unvollkommenheiten bei einem solchen Objectiv statt, die beträchtlich grösser sind als die von der nicht ganz zu hebenden Abweichung wegen der Gestalt herrührenden, so ist es offenbar wichtiger, jene als diese zu berücksichtigen.

Aus dieser Ursache wird es vortheilhafter sein, die Freiheit, die man in der Bestimmung des Verhältnisses der beiden ersten Halbmesser hat, zur Verminderung oder Wegschaffung der Farbenzerstreuung bei den Randstrahlen zu benutzen. In der That hat Hr.\ Prof.\ \textsc{Bohnenberger} durch Rechnung gezeigt, dass in dieser Beziehung das Verhältniss 2 zu 3 dem Verhältnisse 1 zu 3 vorzuziehen ist, indem bei der ersten eine beträchtlich kleinere Farbenzerstreuung der Randstrahlen bewirkt wird, ohne dass die übriggebliebene Abweichung wegen der Kugelgestalt erheblich geworden wäre. Inzwischen bleibt auch bei Hrn.\ Prof.\ \textsc{Bohnenbergers} Einrichtung noch eine Farbenzerstreuung der Randstrahlen zurück,
